%% gallery.tex -- one page, one seed (ConTeXt). One short page per variant, as the
%% LaTeX gallery.tex:
%%   context --florilegiumseed=7 gallery.tex
%% Optional: --galleryink=5B2A1F (the ink, hex), --galleryage=0.8 (the season),
%% --gallerytheme=math (the theme), --galleryoptions=density=0.4 (any setup keys).

\usemodule[florilegium][composition=grown]
\environment env-florilegium

\doifsomething{\getdocumentargument{galleryink}}
  {\definecolor[ink][h=\getdocumentargument{galleryink}]}
\doifsomething{\getdocumentargument{gallerytheme}}
  {\setupflorilegium[theme=\getdocumentargument{gallerytheme}]}
\doifsomething{\getdocumentargument{galleryoptions}}
  {\expanded{\setupflorilegium[\getdocumentargument{galleryoptions}]}}
\setupfootertexts[][]

\starttext
\doifsomething{\getdocumentargument{galleryage}}
  {\florilegiumage{\getdocumentargument{galleryage}}}
\startatsection{3}{1}
\subsection{One seed, one garden}

This page was grown from seed {\tt \doifelsenothing{\getdocumentargument{florilegiumseed}}{page}{\getdocumentargument{florilegiumseed}}}.
Nothing in the margin was placed by hand: three stems rise from a root in the corner and
sway up the column, leaves hang at their nodes, and a greedy loop then finds the largest
remaining hole, chooses a motif that fits it and grows a branch to it from the nearest
stem, until only small gaps are left. The same seed always gives the same garden; a
different seed gives a different one. Because the text on this page ends early, the
garden also grows into the space below these lines.
\florilegiumtextend
\stoptext
